\documentclass[11pt]{article}
\usepackage{mystyle}

\title{Rosen --- \S 1.8: Proof Methods and Strategy\\ \large Selected Exercises}
\author{Alston}
\date{Semester 2, 2026}

\begin{document}
\maketitle

% ======================================================================
\begin{exercise}{3}
    Prove that if $x$ and $y$ are real numbers, then
    $\max(x,y) + \min(x,y) = x + y$.
    [\emph{Hint}: Use a proof by cases, with the two cases corresponding
    to $x \geq y$ and $x < y$, respectively.]
\end{exercise}

% ======================================================================
\begin{exercise}{7}
    Prove the triangle inequality, which states that if $x$ and $y$ are
    real numbers, then $|x| + |y| \geq |x + y|$ (where $|x|$ represents
    the absolute value of $x$, which equals $x$ if $x \geq 0$ and equals
    $-x$ if $x < 0$).
\end{exercise}

% ======================================================================
\begin{exercise}{9}
    Prove that there are 100 consecutive positive integers that are not
    perfect squares. Is your proof constructive or nonconstructive?
\end{exercise}

% ======================================================================
\begin{exercise}{13}
    Prove or disprove that there is a rational number $x$ and an
    irrational number $y$ such that $x^y$ is irrational.
\end{exercise}

% ======================================================================
\begin{exercise}{17}
    Suppose that $a$ and $b$ are odd integers with $a \neq b$. Show
    there is a unique integer $c$ such that $|a - c| = |b - c|$.
\end{exercise}

% ======================================================================
\begin{exercise}{21}
    Prove that given a real number $x$ there exist unique numbers $n$
    and $\ell$ such that $x = n - \ell$, $n$ is an integer, and
    $0 \leq \ell < 1$.
\end{exercise}

% ======================================================================
\begin{exercise}{23}
    The \emph{harmonic mean} of two real numbers $x$ and $y$ equals
    $2xy/(x+y)$. By computing the harmonic and geometric means of
    different pairs of positive real numbers, formulate a conjecture
    about their relative sizes and prove your conjecture.
\end{exercise}

% ======================================================================
\begin{exercise}{25}
    Write the numbers $1, 2, \dots, 2n$ on a blackboard, where $n$ is
    an odd integer. Pick any two of the numbers, $j$ and $k$, write
    $|j - k|$ on the board and erase $j$ and $k$. Continue this process
    until only one integer is written on the board. Prove that this
    integer must be odd.
\end{exercise}

% ======================================================================
\begin{exercise}{29}
    Prove that there is no positive integer $n$ such that
    $n^2 + n^3 = 100$.
\end{exercise}

% ======================================================================
\begin{exercise}{33}
    Adapt the proof in Example 4 in Section 1.7 to prove that if
    $n = abc$, where $a$, $b$, and $c$ are positive integers, then
    $a \leq \sqrt[3]{n}$, $b \leq \sqrt[3]{n}$, or $c \leq \sqrt[3]{n}$.
\end{exercise}

% ======================================================================
\begin{exercise}{35}
    Prove that between every two rational numbers there is an
    irrational number.
\end{exercise}

% ======================================================================
\begin{exercise}{37}
    Let $S = x_1y_1 + x_2y_2 + \cdots + x_ny_n$, where
    $x_1, x_2, \dots, x_n$ and $y_1, y_2, \dots, y_n$ are orderings of
    two different sequences of positive real numbers, each containing
    $n$ elements.
    \begin{enumerate}[label=(\alph*)]
        \item Show that $S$ takes its maximum value over all orderings
        of the two sequences when both sequences are sorted (so that
        the elements in each sequence are in nondecreasing order).
        \item Show that $S$ takes its minimum value over all orderings
        of the two sequences when one sequence is sorted into
        nondecreasing order and the other is sorted into nonincreasing
        order.
    \end{enumerate}
\end{exercise}

% ======================================================================
\begin{exercise}{45}
    Use a proof by exhaustion to show that a tiling using dominoes of a
    $4 \times 4$ checkerboard with opposite corners removed does not
    exist.
    [\emph{Hint}: First show that you can assume that the squares in
    the upper left and lower right corners are removed. Number the
    squares of the original checkerboard from 1 to 16, starting in the
    first row, moving right in this row, then starting in the leftmost
    square in the second row and moving right, and so on. Remove
    squares 1 and 16. To begin the proof, note that square 2 is
    covered either by a domino laid horizontally, which covers squares
    2 and 3, or vertically, which covers squares 2 and 6. Consider
    each of these cases separately, and work through all the subcases
    that arise.]
\end{exercise}

% ======================================================================
\begin{exercise}{47}
    Show that by removing two white squares and two black squares from
    an $8 \times 8$ checkerboard (colored as in the text) you can make
    it impossible to tile the remaining squares using dominoes.
\end{exercise}

% ======================================================================
\begin{exercise}{49}
    \begin{enumerate}[label=(\alph*)]
        \item Draw each of the five different tetrominoes, where a
        tetromino is a polyomino consisting of four squares.
        \item For each of the five different tetrominoes, prove or
        disprove that you can tile a standard checkerboard using these
        tetrominoes.
    \end{enumerate}
\end{exercise}

\end{document}